% =============================================================================
% Преамбула учебника «ПЛИС на практике: Sipeed Tang Nano 20K»
% Компиляция: xelatex main.tex (дважды, для оглавления и ссылок)
% =============================================================================

% --- Языки и шрифты (XeLaTeX) ---
\usepackage{fontspec}
\usepackage{polyglossia}
\setdefaultlanguage{russian}
\setotherlanguage{english}

% На Windows используются Times New Roman / Arial / Consolas,
% на Linux -- метрически совместимые Liberation и DejaVu Sans Mono.
\IfFontExistsTF{Times New Roman}
  {\def\MainFontName{Times New Roman}}{\def\MainFontName{Liberation Serif}}
\IfFontExistsTF{Arial}
  {\def\SansFontName{Arial}}{\def\SansFontName{Liberation Sans}}
\IfFontExistsTF{Consolas}
  {\def\MonoFontName{Consolas}}{\def\MonoFontName{DejaVu Sans Mono}}
\setmainfont{\MainFontName}[Ligatures=TeX]
\setsansfont{\SansFontName}[Ligatures=TeX]
\setmonofont{\MonoFontName}[Scale=MatchLowercase]
% polyglossia требует явно указать шрифты для кириллицы
\newfontfamily\cyrillicfont{\MainFontName}[Ligatures=TeX]
\newfontfamily\cyrillicfontsf{\SansFontName}[Ligatures=TeX]
\newfontfamily\cyrillicfonttt{\MonoFontName}[Scale=MatchLowercase]

\usepackage{microtype}
\emergencystretch=3em

% --- Абзацы ---
\usepackage{indentfirst}
\usepackage{setspace}
\setstretch{1.2}
\setlength{\parindent}{1.25cm}
\setlength{\parskip}{3pt plus 1pt}

% --- Математика ---
\usepackage{amsmath}
\usepackage{amssymb}

% --- Поля ---
\usepackage{geometry}
\geometry{a4paper, top=22mm, bottom=22mm, left=25mm, right=20mm,
          headheight=15mm, footskip=12mm}

% --- Цвета ---
\usepackage[table]{xcolor}
\definecolor{FpgaBlue}{RGB}{0, 91, 187}
\definecolor{FpgaDark}{RGB}{22, 43, 58}
\definecolor{FpgaOrange}{RGB}{255, 140, 0}
\definecolor{FpgaGreen}{RGB}{0, 140, 70}
\definecolor{FpgaRed}{RGB}{200, 30, 45}
\definecolor{FpgaPurple}{RGB}{120, 40, 160}
\definecolor{FpgaLight}{RGB}{242, 246, 252}
\definecolor{FpgaGray}{RGB}{120, 120, 125}
\definecolor{BoardGreen}{RGB}{20, 90, 60}
\definecolor{ChipBlack}{RGB}{35, 35, 40}
\definecolor{CodeBack}{RGB}{248, 249, 251}
\definecolor{CodeComment}{RGB}{0, 128, 0}
\definecolor{CodeKeyword}{RGB}{0, 60, 170}
\definecolor{CodeString}{RGB}{170, 60, 0}

% --- Графика ---
\usepackage{etoolbox}
\usepackage{xstring}
\usepackage{graphicx}
\usepackage{float}
\usepackage{tikz}
\usetikzlibrary{arrows.meta, calc, positioning, shapes.geometric, shapes.misc,
  shapes.gates.logic.US, shapes.gates.logic.IEC, shapes.symbols,
  fit, backgrounds, decorations.pathreplacing, patterns, shadows, matrix}
\usepackage{tikz-timing}
\usepackage{pgfplots}
\pgfplotsset{compat=1.18}
\usepgfplotslibrary{groupplots}
\usepackage[siunitx, RPvoltages]{circuitikz}

% --- Заголовки, колонтитулы ---
\usepackage[automark, headsepline]{scrlayer-scrpage}
\clearpairofpagestyles
\ihead{\headmark}
\ohead{\textbf{\textcolor{FpgaBlue}{TANG NANO 20K}}}
\cfoot{\pagemark}
\setkomafont{pageheadfoot}{\sffamily\small\color{FpgaDark}}
\setkomafont{pagenumber}{\sffamily\bfseries\color{FpgaBlue}}
\setkomafont{chapter}{\sffamily\huge\bfseries\color{FpgaBlue}}
\setkomafont{section}{\sffamily\Large\bfseries\color{FpgaBlue}}
\setkomafont{subsection}{\sffamily\large\bfseries\color{FpgaDark}}
\setkomafont{subsubsection}{\sffamily\normalsize\bfseries\color{FpgaOrange}}
\setkomafont{paragraph}{\sffamily\bfseries\color{FpgaDark}}
\setkomafont{part}{\sffamily\Huge\bfseries\color{FpgaBlue}}
\renewcommand*{\chapterformat}{%
  \colorbox{FpgaBlue}{\color{white}\,\thechapter\,}\enskip}

% --- Списки и таблицы ---
\usepackage{enumitem}
\setlist{itemsep=1pt, parsep=0pt, topsep=3pt}
\setlist[itemize]{label=\textcolor{FpgaBlue}{\textbullet}, leftmargin=*}
\setlist[enumerate]{label=\textcolor{FpgaBlue}{\bfseries\arabic*.}, leftmargin=*}
\setlist[description]{font=\bfseries\color{FpgaDark}}
\usepackage{booktabs}
\usepackage{tabularx}
\usepackage{multirow}
\usepackage{longtable}
\usepackage{caption}
\captionsetup{font=small, labelfont={bf,color=FpgaBlue}}
\usepackage{subcaption}

% --- Листинги кода (minted: нужен Pygments и ключ -shell-escape) ---
\usepackage[outputdir=../build]{minted}
\usemintedstyle{friendly}
\setminted{fontsize=\small, breaklines=true, tabsize=4, baselinestretch=1.0}

% --- Цветные блоки ---
\usepackage[most]{tcolorbox}
\tcbuselibrary{minted}
\tcbset{codestyle/.style={enhanced, breakable, listing only, listing engine=minted,
  colback=CodeBack, colframe=FpgaBlue, boxrule=0pt, leftrule=2.2pt, arc=0pt,
  left=6pt, right=4pt, top=2pt, bottom=2pt, before skip=6pt, after skip=8pt,
  fonttitle=\ttfamily\small\bfseries, coltitle=FpgaDark, colbacktitle=FpgaLight,
  attach boxed title to top left={yshift=-1pt}, boxed title style={boxrule=0pt, arc=0pt}}}
% Фрагмент кода: \begin{vcode} ... \end{vcode}, язык по умолчанию Verilog
\newtcblisting{vcode}[1][verilog]{codestyle, minted language=#1,
  minted options={fontsize=\small, breaklines=true, tabsize=4}}
% Файл целиком с номерами строк: \vfile{путь}{подпись}
\newtcbinputlisting{\vfile}[2]{codestyle, minted language=verilog,
  minted options={fontsize=\small, breaklines=true, tabsize=4, linenos, numbersep=5pt},
  listing file={#1}, title={#2}, left=20pt}
\newtcolorbox{infobox}[1][Справка]{
  enhanced, breakable, colback=FpgaLight, colframe=FpgaBlue,
  fonttitle=\sffamily\bfseries, title={#1}, arc=2mm, boxrule=0.8pt}
\newtcolorbox{warnbox}[1][Внимание!]{
  enhanced, breakable, colback=orange!6, colframe=FpgaOrange,
  fonttitle=\sffamily\bfseries, title={#1}, arc=2mm, boxrule=0.8pt}
\newtcolorbox{taskbox}[1][Задание]{
  enhanced, breakable, colback=green!4, colframe=FpgaGreen,
  fonttitle=\sffamily\bfseries, title={#1}, arc=2mm, boxrule=0.8pt}
\newtcolorbox{ideabox}[1][Главная мысль]{
  enhanced, breakable, colback=purple!4, colframe=FpgaPurple,
  fonttitle=\sffamily\bfseries, title={#1}, arc=2mm, boxrule=0.8pt}

% --- Прочее ---
\usepackage{siunitx}
\sisetup{output-decimal-marker={,}, group-separator={\,}, locale=FR}
\usepackage{hyperref}
\hypersetup{colorlinks=true, linkcolor=FpgaBlue, urlcolor=FpgaBlue,
  citecolor=FpgaBlue, pdftitle={ПЛИС на практике: Sipeed Tang Nano 20K},
  pdfauthor={Корягин С.А.}}

% Удобные команды
\newcommand{\code}[1]{\texttt{\small #1}}
\newcommand{\term}[1]{\textbf{\textcolor{FpgaDark}{#1}}}
\newcommand{\codepath}{../code}

% Общие стили TikZ
\tikzset{
  block/.style={draw=FpgaBlue, fill=FpgaLight, thick, rounded corners=2pt,
                minimum height=9mm, minimum width=22mm, align=center,
                font=\sffamily\small},
  hblock/.style={block, fill=FpgaBlue, text=white},
  oblock/.style={block, draw=FpgaOrange, fill=orange!10},
  gblock/.style={block, draw=FpgaGreen, fill=green!8},
  arr/.style={-{Stealth[length=2.5mm]}, thick, color=FpgaDark},
  wire/.style={thick, color=FpgaDark},
  bus/.style={line width=2pt, color=FpgaBlue},
  lbl/.style={font=\sffamily\scriptsize, color=FpgaDark},
}
\tikzset{
  gbase/.style={draw=FpgaDark, thick, fill=FpgaLight, logic gate inputs=nn,
                minimum width=12mm, minimum height=9mm, inner sep=0pt},
  gAND/.style={and gate US, gbase},
  gOR/.style={or gate US, gbase},
  gXOR/.style={xor gate US, gbase},
  gNAND/.style={nand gate US, gbase},
  gNOR/.style={nor gate US, gbase},
  gXNOR/.style={xnor gate US, gbase},
  ibase/.style={draw=FpgaDark, thick, fill=FpgaLight, logic gate inputs=nn,
                minimum width=8mm, minimum height=10mm, font=\sffamily\small},
  iAND/.style={and gate IEC, ibase},
  iOR/.style={or gate IEC, ibase},
  iXOR/.style={xor gate IEC, ibase},
  iNAND/.style={nand gate IEC, ibase},
  iNOR/.style={nor gate IEC, ibase},
  iXNOR/.style={xnor gate IEC, ibase},
  iNOT/.style={not gate IEC, draw=FpgaDark, thick, fill=FpgaLight,
               minimum width=8mm, minimum height=7mm, font=\sffamily\small},
  gNOT/.style={not gate US, draw=FpgaDark, thick, fill=FpgaLight,
               minimum width=9mm, minimum height=6mm, inner sep=0pt},
}
